% -*- coding: utf-8 -*-
% 南开大学本科毕业论文 LaTeX 模板；使用 XeLaTeX + Biber 编译。
\documentclass[12pt,openright]{book}

\usepackage{ifxetex}
\ifxetex
  \usepackage[bookmarksnumbered]{hyperref}
\else
  \usepackage[unicode,bookmarksnumbered]{hyperref}
\fi
% 上下2.54cm；左右3.17cm；页眉1.5cm；页脚1.75cm
\usepackage{geometry}
\usepackage[emptydoublepage]{NKUThesis}           % 中文
% \usepackage[emptydoublepage,English]{NKUThesis} % 英文
% 参考文献沿用模板的 NKUThesis 样式，按正文中的引用顺序排列。
\usepackage[maxnames=3,minnames=3,sorting=none,style=NKUThesis]{biblatex}

\hypersetup{
  colorlinks=true,
  pdfborder=0 0 1,
  citecolor=black,
  linkcolor=black,
  pdftitle={南开大学本科毕业论文 LaTeX 模板使用说明},
  pdfauthor={},
  pdfsubject={本科毕业论文模板使用说明},
  pdfkeywords={南开大学,本科毕业论文,LaTeX模板}
}

\usepackage{amsmath}
\usepackage{amssymb}
\usepackage{booktabs}
\usepackage{graphicx}
\usepackage{multirow}
\usepackage{pgfplots}
\pgfplotsset{compat=1.13}
\usepackage[labelsep=space]{caption}
\DeclareCaptionFont{wuhao}{\fontsize{10.5}{12}\selectfont}
\captionsetup{font=wuhao}
\captionsetup[figure]{justification=centering}
\usepackage{subfigure}
\usepackage{color,xcolor}
\usepackage{listings}
% 普通程序代码：五号字、单倍行距，带边框、行号和语法高亮。
% 使用 lstlisting 的 language 选项指定语言，如 Python、C++、bash。
\definecolor{NKTCodeBackground}{HTML}{F7F8FA}
\definecolor{NKTCodeBorder}{HTML}{CBD2D9}
\definecolor{NKTCodeKeyword}{HTML}{1F4E79}
\definecolor{NKTCodeComment}{HTML}{50734C}
\definecolor{NKTCodeString}{HTML}{9C3D10}
\definecolor{NKTCodeNumber}{HTML}{666666}
\lstset{
  basicstyle=\linespread{1}\fontsize{10.5}{12}\selectfont\ttfamily,
  lineskip=0pt,
  keywordstyle=\color{NKTCodeKeyword}\bfseries,
  commentstyle=\color{NKTCodeComment}\itshape,
  stringstyle=\color{NKTCodeString},
  backgroundcolor=\color{NKTCodeBackground},
  frame=single,
  rulecolor=\color{NKTCodeBorder},
  framerule=0.4pt,
  framesep=5pt,
  xleftmargin=24pt,
  xrightmargin=6pt,
  framexleftmargin=18pt,
  numbers=left,
  numberstyle=\normalfont\fontsize{8}{10}\selectfont\color{NKTCodeNumber},
  numbersep=7pt,
  stepnumber=1,
  firstnumber=1,
  tabsize=4,
  columns=fullflexible,
  keepspaces=true,
  showstringspaces=false,
  breaklines=true,
  breakatwhitespace=false,
  aboveskip=8pt,
  belowskip=8pt,
}
\usepackage{tikz}
\usepackage{setspace}
% 算法伪代码沿用正文小四号；普通程序代码使用 lstlisting。
\usepackage{algorithm}
\usepackage{algpseudocode}
\usepackage{cases}
\usepackage{tabularx}

% 绘图配色与图案，可供正文中的图表使用。
\definecolor{corange}{HTML}{FF9666}
\definecolor{lblue}{HTML}{97FFFF}
\definecolor{cblue}{HTML}{AB82FF}
\definecolor{cred}{HTML}{FF99A1}
\definecolor{cgreen}{HTML}{7ACCBE}
\definecolor{red1}{RGB}{254,67,101}
\definecolor{red2}{RGB}{252,157,154}
\definecolor{red3}{RGB}{249,205,173}
\definecolor{grey1}{RGB}{200,200,169}
\definecolor{grey2}{RGB}{131,175,155}
\definecolor{color1}{HTML}{0099CC}
\definecolor{color2}{HTML}{F9B46F}
\definecolor{cblue1}{HTML}{0074D9}
\definecolor{lblue1}{HTML}{99CCFF}
\definecolor{corange1}{HTML}{FFDDCE}
\definecolor{corange2}{HTML}{FF9900}
\definecolor{corange3}{HTML}{FF9966}
\definecolor{corange4}{HTML}{FF6600}
\definecolor{cred2}{HTML}{FF6666}
\usetikzlibrary{patterns}

\addbibresource{references.bib}
% 参考文献表只列出正文实际引用的条目。
\DeclareBibliographyCategory{cited}
\AtEveryCitekey{\addtocategory{cited}{\thefield{entrykey}}}

% 按章编号的定理类环境；名称和编号的字体由 NKUThesis.sty 设置。
\newtheorem{Theorem}{\hskip 2em 定理}[chapter]
\newtheorem{Lemma}[Theorem]{\hskip 2em 引理}
\newtheorem{Corollary}[Theorem]{\hskip 2em 推论}
\newtheorem{Proposition}[Theorem]{\hskip 2em 命题}
\newtheorem{Definition}[Theorem]{\hskip 2em 定义}
\newtheorem{Example}[Theorem]{\hskip 2em 例}
\newcommand{\upcite}[1]{\textsuperscript{\textsuperscript{\cite{#1}}}}
\renewcommand{\algorithmicrequire}{\textbf{输入：}}
\renewcommand{\algorithmicensure}{\textbf{输出：}}
\algnewcommand{\LeftComment}[1]{\Statex \(\triangleright\) #1}
\floatname{algorithm}{算法}
\newcommand{\innersection}[2]{\noindent\textbf{#1. #2}\par}

% 提示词示例框；其中的正文及 verbatim 片段沿用小四号。
\usepackage[many,breakable,skins]{tcolorbox}
\newcounter{myboxcounter}[chapter]
\renewcommand{\themyboxcounter}{\thechapter.\arabic{myboxcounter}}
\tcbset{
  promptbox_base/.style={
    colback=gray!10,
    colframe=gray!50!black,
    boxrule=1.0pt,
    arc=2mm,
    left=5pt,
    right=5pt,
    top=5pt,
    bottom=5pt,
    fonttitle=\bfseries,
    enhanced,
  }
}
\makeatletter
\newenvironment{promptbox}[2][]{%
  \refstepcounter{myboxcounter}%
  \edef\@currentlabel{\themyboxcounter}%
  \begin{tcolorbox}[
    promptbox_base,
    breakable, % 允许跨页
    title={例 \themyboxcounter: #2},
    title after break={例 \themyboxcounter: #2（续）}, % 添加续标题
    #1
  ]%
}{%
  \end{tcolorbox}%
}
\makeatother

\usepackage{pifont}

% 自定义颜色
\definecolor{safeGreen}{rgb}{0.0,0.5,0.0}    % 深绿色
\definecolor{unsafeRed}{rgb}{0.8,0.0,0.0}    % 深红色
\definecolor{niceOrange}{rgb}{0.95,0.5,0.15} % 橙色
\definecolor{niceBlue}{rgb}{0.0,0.35,0.7}    % 蓝色
\usepackage{hhline}

\begin{document}

% 设置封面基本信息：以下均为占位内容，请按实际情况填写。
% 注意：英文逗号“,”是字段分隔符。字段内容含英文逗号时，
% 请用花括号保护该逗号（如 {,}），或包围整个字段值（如 {Title, Subtitle}）。
\NKTsetup{
  % 中英文题目可手动分成两行；不需要第二行时填“无”，该行将不显示。
  论文题目(中文) = 南开大学本科毕业论文模板使用说明,
  论文题目(中文)(第二行) = 无,
  论文题目(英文) = Nankai University Undergraduate Thesis Template,
  论文题目(英文)(第二行) = User Guide,
  学号 = 0000000,
  姓名 = 姓名,
  % 年级填写入学年份（4位数字）；“级”为显示后缀，一般保留。
  年级 = 0000,
  级 = 级,
  专业 = 专业名称,
  系别 = 系别名称,
  学院 = 学院名称,
  % 导师按“姓名\quad 职称”填写；\quad 用于在姓名与职称间留空。
  指导教师 = 教师姓名\quad 职称,
  % 有校外导师时填写，没有则填“无”。
  % 示例：校外导师 = 教师姓名\quad 职称（单位名称）,
  校外导师 = 无,
  % “论文完成时间”填写年份，“月份”填写月份数字；“年”“月”为显示后缀。
  % 若已在月份值中写入“月”字，则应将“月”字段留空，避免重复显示。
  论文完成时间 = 0000,
  年 = 年,
  月份 = 5,
  月 = 月,
  % 无后缀的信息栏共用此空值，请保留为空。
  空后缀 =
}

% -*- coding: utf-8 -*-
% 将以下示例替换为论文的中英文摘要和关键词。
\begin{zhaiyao}
\begin{spacing}{1.5}
本模板用于南开大学本科毕业论文的排版，面向计算机学院和密码与网络空间安全学院的格式要求进行定制。本文仅提供简短的使用说明，演示正文、图表、公式、代码、参考文献及附录的组织方式。正式撰写时，请替换封面信息、摘要、正文和致谢，并对照所在学院的最新要求检查格式。
\end{spacing}
\end{zhaiyao}

\begin{guanjianci}
本科毕业论文；模板；排版
\end{guanjianci}

\begin{abstract}
\begin{spacing}{1.5}
This template supports undergraduate thesis typesetting at Nankai University.
It provides a starting point for thesis writing and should be adapted to the requirements of the relevant college.
This short guide demonstrates the organization of the main text, figures, tables, equations, code, references, and appendices.
Before writing a thesis, replace the sample content and check the current requirements of the relevant college.
\end{spacing}
\end{abstract}

\begin{keywords}
undergraduate thesis; template; typesetting
\end{keywords}

\tableofcontents

\begin{spacing}{1.5}
% -*- coding: utf-8 -*-
% 正文从这里开始；可直接替换本文件，也可在 tex/ 中继续拆分章节。
\chapter{模板使用说明}
\label{chap:guide}

\section{适用范围}
本模板面向南开大学\textbf{计算机学院}和\textbf{密码与网络空间安全学院}的本科毕业论文，主要参考《南开大学本科毕业论文（设计）指导手册（2025）》\upcite{nku-guide-2025}以及《本科生论文格式要求补充规定2023V2》\upcite{nku-supplement-2023}。其他学院使用时，可能需要按本学院要求修改。

\section{文件与编译}
在 \path{main.tex} 中填写封面信息，在 \path{tex/abstract.tex} 中填写摘要，在本文件中撰写正文。插图放入 \path{figures/}，文献条目写入 \path{references.bib}。使用 XeLaTeX 和 Biber 编译，具体命令见项目 README 和附录。

\subsection{标题与交叉引用}
使用 \verb|\chapter|、\verb|\section| 和 \verb|\subsection| 组织三级标题。为图、表、公式添加 \verb|\label|，再用 \verb|\ref| 引用，避免手工填写编号。

\subsection{中英文文献引用}
文献使用 \verb|\cite| （行内引用\cite{BERT}） 或上标引用命令 \verb|\upcite| （上标引用\upcite{BERT}）。推荐使用 \verb|\upcite{条目键}| 作上标引用。中文文献示例：赵月等总结了大语言模型的安全现状与挑战\upcite{ZW-AQTZ}。英文文献示例：Devlin 等提出了 BERT 预训练方法\upcite{BERT}。

上述条目位于 \path{references.bib}，对应的条目键为 \verb|ZW-AQTZ| 和 \verb|BERT|。同一处引用多篇文献时，在 \verb|\upcite| 中用英文逗号分隔条目键，多条文献引用显示效果\upcite{Transformer,LoRA,DeepSeek-R1}。

\chapter{基本排版示例}

\section{公式与定义}
下面给出均值公式的排版示例：
\begin{equation}
  \bar{x}=\frac{1}{n}\sum_{i=1}^{n}x_i.
  \label{eq:mean}
\end{equation}
式~\ref{eq:mean} 的编号随章节自动生成。定义、定理和引理可分别使用 \verb|Definition|、\verb|Theorem| 和 \verb|Lemma| 环境。
\begin{Definition}
样本均值是所有样本数值之和与样本数量之比。
\end{Definition}

\section{图表}
图~\ref{fig:placeholder} 是插图占位示例。实际插图时，将方框替换为 \verb|\includegraphics| 命令，并在图下方填写图题。
\begin{figure}[htbp]
  \centering
  % 例如：\includegraphics[width=0.6\linewidth]{figures/example.pdf}
  \fbox{\parbox[c][16mm][c]{0.6\linewidth}{\centering 插图占位示意}}
  \caption{插图示例}
  \label{fig:placeholder}
\end{figure}

表~\ref{tab:example} 演示三线表。表题放在表格上方，表格内容的字号由模板设置。
\begin{table}[htbp]
  \centering
  \caption{表格示例}
  \label{tab:example}
  \begin{tabular}{lc}
    \toprule
    项目 & 示例值 \\
    \midrule
    样本数量 & 3 \\
    样本均值 & 2 \\
    \bottomrule
  \end{tabular}
\end{table}

\section{程序代码}
普通程序代码使用 \verb|lstlisting|，保持五号字、单倍行距，并显示边框、行号和语法高亮。通过 \verb|language=Python| 等选项指定语言，例如：

\noindent\begin{minipage}{\linewidth}
\begin{lstlisting}[language=Python]
# Calculate the mean of a nonempty list.
def mean(values):
    return sum(values) / len(values)

print("Mean:", mean([1, 2, 3]))
\end{lstlisting}
\end{minipage}


% 附录位于正文之后、参考文献之前；无附录时可注释下一行。
% -*- coding: utf-8 -*-
% 本模板的附录命令直接接收标题，与标准 book 的用法不同。
\appendix{补充说明}
\label{app:build}

附录可用于放置补充推导、较长代码或其他不适合放入正文的材料。使用 \verb|\appendix{...}| 设置附录标题；如需多个附录，可继续使用此命令。

在项目根目录运行以下命令即可编译：
\begin{lstlisting}[language=bash]
latexmk -xelatex main.tex
\end{lstlisting}

未使用 latexmk 时，依次运行：
\begin{lstlisting}[language=bash]
xelatex main.tex
biber main
xelatex main.tex
xelatex main.tex
\end{lstlisting}

编译结果为 \path{main.pdf}。不需要附录时，注释 \path{main.tex} 中的附录引入语句即可。


% 参考文献
\clearpage
\phantomsection
\zihaowu
\printbibliography[category=cited]

% -*- coding: utf-8 -*-
\begin{zhixie}
\begin{spacing}{1.5}
请在此填写致谢，简要说明在论文研究与写作过程中获得的指导、支持和帮助。正式提交前，请将本段说明替换为自己的致谢内容。
\end{spacing}
\end{zhixie}


\end{spacing}

\end{document}
